% This document was generated by an LLM.

\documentclass{essay}

\title{Undoing the Derivative}
\subtitle{What Becomes of the Fundamental Theorem in Several Variables}
\author{}
\date{}

\begin{document}

\maketitle

\section{Introduction: A Deceptively Simple Question}

In single-variable calculus, differentiation and integration are inverse operations. If we differentiate a function and then integrate the result, we recover the original function; the only information lost along the way is an additive constant. This is the content of the Fundamental Theorem of Calculus (FTC), and it is so familiar that it is easy to regard it as obvious.

The question this essay pursues is what becomes of this inverse relationship when one variable becomes two. Given a function $f(x,y)$, we can differentiate it in several inequivalent ways: we can take a single partial derivative $\partial f/\partial x$; we can take the gradient $\nabla f = (f_x, f_y)$; we can iterate partials, as in the mixed derivative $\partial^2 f/\partial x\,\partial y$. Which of these operations can be undone? What is the analogue of the indefinite integral? What plays the role of the constant of integration? And do the ``indefinite multiple integrals'' conspicuously absent from most textbooks exist at all?

The answer, developed over the following sections, is that the tidy one-variable picture does not generalize as a single fact. It fragments into several distinct inverse problems---one for each differential operator---each with its own notion of antiderivative, its own version of the fundamental theorem, and its own description of the lost information. Remarkably, the fragments are then reunified at a higher level of abstraction, where a single operator (the exterior derivative) and a single theorem (the generalized Stokes theorem) subsume every case.

Throughout, one lens organizes everything:
\begin{quote}
\emph{The information lost when applying a differential operator is precisely the \textbf{kernel} of that operator---the set of objects it sends to zero. To understand how to invert an operator, first understand what it kills.}
\end{quote}
In one variable the kernel of $d/dx$ is the constants, whence ``$+C$.'' In several variables the kernels will grow, shrink, and, in the final act, acquire topological meaning.

\section{The One-Variable Case, Reread Carefully}

Before generalizing a theorem, one should state precisely what is being generalized. The FTC is really two statements, and it pays to separate them.

\begin{theorem}[FTC, first half]
If $f$ is continuous on an interval containing $a$, then
\[
\frac{d}{dx}\int_a^x f(t)\dd t = f(x).
\]
\end{theorem}

\begin{theorem}[FTC, second half]
If $F$ is continuously differentiable, then
\[
\int_a^x F'(t)\dd t = F(x) - F(a).
\]
\end{theorem}

The first half says that differentiation undoes integration \emph{exactly}: no information is lost in that direction. The second half says that integration undoes differentiation \emph{up to the value $F(a)$}---a single number. If we only know $F'$, we know $F$ up to an additive constant.

Why a constant, exactly? Because two antiderivatives of the same function differ by a function whose derivative vanishes identically, and (on an interval) the only such functions are constants. In the language of linear algebra: the operator $d/dx$ is linear, and its \emph{kernel}---the solution set of $F' = 0$---consists of the constant functions. The familiar ``$+C$'' is not a bookkeeping ritual; it is a parametrization of this kernel.

Two further rereadings will matter later.

\emph{First}, the indefinite integral is best understood as a definite integral with a variable upper endpoint. The notation $\int f(x)\dd x$ suggests a formal symbol-pushing operation, but the honest construction of an antiderivative is
\[
F(x) = \int_a^x f(t)\dd t,
\]
a genuine function of $x$ defined by accumulation from a base point $a$. Different choices of base point change $F$ by a constant---the kernel again. This ``variable endpoint'' viewpoint is the one that generalizes.

\emph{Second}, note the small print ``on an interval.'' If the domain of $F$ is a union of two disjoint intervals, then $F'=0$ allows a different constant on each piece, and the kernel is two-dimensional. Already in one variable, the lost information is sensitive to the shape of the domain. This innocuous remark is the seed of Section~\ref{sec:topology}.

\section{Inverting One Partial at a Time}

The gentlest step into two variables is to differentiate with respect to one variable only. Fix the operator $\partial/\partial x$ acting on functions $f(x,y)$ defined on, say, a rectangle or all of $\R^2$. Can it be inverted?

Yes, and by the obvious candidate: integrate with respect to $x$ while treating $y$ as a fixed parameter. If
\[
\frac{\partial f}{\partial x} = P(x,y),
\]
then for each fixed $y$ the one-variable FTC applies along the horizontal line at height $y$, giving
\[
f(x,y) = \int_{a}^{x} P(s,y)\dd s + (\text{something independent of } x).
\]
What is the ``something''? It is anything killed by $\partial/\partial x$: any function of $y$ alone. Hence the general solution is
\[
f(x,y) = \int P(x,y)\dd x + g(y), \qquad g \text{ arbitrary}.
\]

This \emph{is} an indefinite integral in the multivariable setting; it simply does not receive that name in textbooks. And the structure of the answer is the first genuinely new phenomenon: \textbf{the constant of integration has inflated into an arbitrary function}. The kernel of $\partial/\partial x$ is the infinite-dimensional space of functions of $y$, vastly larger than the one-dimensional space of constants.

This inflation is not a curiosity; it is the structural reason why general solutions of partial differential equations contain arbitrary \emph{functions} where general solutions of ordinary differential equations contain arbitrary \emph{constants}. Solving a PDE means inverting a partial differential operator, and each inversion contributes its kernel to the general solution. The simplest PDE of all, $u_x = 0$, has general solution $u = g(y)$: the kernel itself.

\section{The Mixed Partial and the True Indefinite Double Integral}
\label{sec:mixed}

Now iterate. Consider the mixed partial operator
\[
L = \frac{\partial^2}{\partial x\,\partial y},
\]
and the problem: given $h(x,y)$, find all $f$ with $f_{xy} = h$ (writing $f_{xy}$ for $\partial^2 f/\partial x\,\partial y$).

Since $L$ is a composition of two partial derivatives, its inverse is a composition of two partial antiderivatives---an \emph{iterated indefinite integral}. Integrate first in $y$, holding $x$ fixed:
\[
\frac{\partial f}{\partial x} = \int h(x,y)\dd y + \varphi(x),
\]
where $\varphi$ is an arbitrary function of $x$ (the kernel of $\partial/\partial y$). Now integrate in $x$, holding $y$ fixed:
\[
f(x,y) = \iint h(x,y)\dd y \dd x \;+\; g_1(x) \;+\; g_2(y),
\]
where $g_1(x) = \int \varphi(x)\dd x$ and $g_2(y)$ is the arbitrary function contributed by the second integration. (An arbitrary function of $x$, integrated in $x$, is just another arbitrary function of $x$; so no generality is lost in writing the answer this way.)

This is the closest thing in the subject to a true ``indefinite double integral,'' and the lost information is now
\[
\ker L = \{\, g_1(x) + g_2(y) \,\}:
\]
\emph{two} arbitrary one-variable functions. The verification is immediate: $\partial/\partial y$ kills $g_1(x)$, and $\partial/\partial x$ kills $g_2(y)$, so $L$ kills their sum. Conversely, if $f_{xy}=0$, then $f_x$ is independent of $y$, say $f_x = \varphi(x)$, and integrating gives $f = \int\varphi + g_2(y)$, which has the stated form.

\begin{example}
Solve $f_{xy} = 6xy$. Integrating in $y$: $f_x = 3xy^2 + \varphi(x)$. Integrating in $x$:
\[
f(x,y) = \tfrac{3}{2}x^2y^2 + g_1(x) + g_2(y).
\]
One checks directly that $f_{xy} = 6xy$ regardless of the choices of $g_1, g_2$.
\end{example}

\subsection*{A two-dimensional fundamental theorem}

There is also a definite-integral counterpart, a genuine 2D analogue of the FTC's first half. Define, with variable upper limits,
\[
F(x,y) = \int_a^x \! \int_b^y h(s,t)\dd t \dd s.
\]
Then differentiating once in $x$ (by the one-variable FTC, applied to the outer integral) and once in $y$ (applied to the inner one) gives
\[
\frac{\partial^2 F}{\partial x\,\partial y} = h(x,y).
\]
The definite double integrals over rectangles that populate every multivariable textbook secretly contain this inverse relationship. It goes unremarked because double integrals are introduced for a different purpose---measuring volume, mass, or probability over a region---and their role as antidifferentiation is incidental to that agenda. But the relationship is there: the double integral with variable corners is the indefinite integral of the mixed partial.

\subsection*{d'Alembert's formula as a kernel computation}

The kernel $\{g_1(x)+g_2(y)\}$ is not an idle formula; it is the solution of the wave equation in disguise. The one-dimensional wave equation $u_{tt} = c^2 u_{xx}$ becomes, in the characteristic coordinates $\xi = x - ct$, $\eta = x + ct$, simply
\[
u_{\xi\eta} = 0.
\]
By the computation above, the general solution is $u = F(\xi) + G(\eta)$, that is,
\[
u(x,t) = F(x - ct) + G(x + ct):
\]
a superposition of two waves traveling in opposite directions at speed $c$. D'Alembert's celebrated formula is, at bottom, the statement that the kernel of the mixed partial consists of sums of one-variable functions. The ``arbitrary functions'' of PDE theory and the ``constants of integration'' of ODE theory are one and the same phenomenon, seen at different scales.

\section{The Gradient: An Overdetermined Inverse Problem}
\label{sec:gradient}

The operators considered so far take one function to one function. The gradient is different in kind: it takes one function to a \emph{vector} of functions,
\[
\nabla f = \left( \frac{\partial f}{\partial x}, \frac{\partial f}{\partial y} \right),
\]
and this changes the character of the inverse problem entirely. Given a vector field $(P,Q)$, the equation $\nabla f = (P,Q)$ is a system of two equations in one unknown function: it is \emph{overdetermined}, and we should expect it to be solvable only when $P$ and $Q$ satisfy a compatibility condition. When it is solvable, $(P,Q)$ is called a \emph{gradient field} (or conservative field) and $f$ a \emph{potential} for it.

\subsection{The recovery algorithm}

Suppose first that a potential exists. It can be found by iterated partial integration, with a matching step.

\begin{example}
Let $\nabla f = (2xy + 3x^2,\; x^2 + \cos y)$. Integrate the first component in $x$:
\[
f(x,y) = x^2 y + x^3 + g(y),
\]
with $g$ the unavoidable arbitrary function. Now impose the second equation: differentiating in $y$,
\[
f_y = x^2 + g'(y) \stackrel{!}{=} x^2 + \cos y \quad\Longrightarrow\quad g'(y) = \cos y,
\]
so $g(y) = \sin y + C$ and
\[
f(x,y) = x^2 y + x^3 + \sin y + C.
\]
\end{example}

Notice what had to happen at the matching step: after substituting the partial result into the second equation, every occurrence of $x$ cancelled, leaving an equation for $g'(y)$ in terms of $y$ alone. Had a stray $x$ survived, the procedure would have collapsed---no function $g(y)$ could satisfy an equation involving $x$. The success of the cancellation is not luck; it is the visible trace of a hidden compatibility condition.

Notice also the size of the lost information: a single constant $C$, exactly as in one variable. The kernel of $\nabla$ on a connected domain is just the constants, since $\nabla f = 0$ forces $f$ to be constant along every path. (On a disconnected domain: one constant per connected component, echoing the one-variable small print.)

\subsection{The obstruction}

What is the hidden condition? If $f$ is twice continuously differentiable, Clairaut's theorem says its mixed partials commute:
\[
\frac{\partial^2 f}{\partial y\,\partial x} = \frac{\partial^2 f}{\partial x\,\partial y}.
\]
Applying this to $\nabla f = (P,Q)$, where $P = f_x$ and $Q = f_y$, yields the necessary condition
\[
\boxed{\;\frac{\partial P}{\partial y} = \frac{\partial Q}{\partial x}.\;}
\]
If this fails anywhere, no potential exists---full stop.

\begin{example}
The field $(P,Q) = (-y, x)$, which rotates rigidly about the origin, has $P_y = -1$ and $Q_x = 1$. It is not the gradient of anything: \emph{gradients cannot swirl}. Geometrically, a gradient field always points in the direction of steepest increase of some landscape $f$; following it can never return you to your starting point at the same altitude, whereas the rotation field takes you in circles.
\end{example}

The contrast with Section~\ref{sec:mixed} deserves emphasis. \emph{Every} continuous $h$ is the mixed partial of something---the inverse problem for $L$ is always solvable, and its kernel is huge. For the gradient the situation is reversed: \emph{most} vector fields are not gradients---the inverse problem is only conditionally solvable---but when it is solvable, the kernel is tiny. Growing kernels and shrinking solvability are two sides of the same accounting.

\subsection{Path integrals and path independence}

The gradient has its own fundamental theorem, and its own ``definite integral with a variable endpoint.'' The \emph{gradient theorem} (the FTC for line integrals) states that for any smooth curve $\gamma$ from a point $A$ to a point $B$,
\[
\int_\gamma \nabla f \cdot d\mathbf{r} \;=\; f(B) - f(A).
\]
The right-hand side depends only on the endpoints, not the route: line integrals of gradient fields are \emph{path independent}. Conversely, if the line integrals of a continuous field $(P,Q)$ are path independent, one may fix a base point $(x_0,y_0)$ and \emph{define}
\[
f(x,y) = \int_{(x_0,y_0)}^{(x,y)} P\dd x + Q\dd y
\]
along any path; the result is a well-defined potential. This is the exact analogue of constructing an antiderivative as $\int_a^x f(t)\dd t$: a definite integral from a base point to a variable endpoint, with the choice of base point accounting for the constant $C$.

Path independence gives a second diagnosis of the rotation field: integrating $(-y,x)$ counterclockwise around the unit circle gives
\[
\oint (-y)\dd x + x \dd y = \int_0^{2\pi} \big( \sin^2\theta + \cos^2\theta \big)\dd\theta = 2\pi \neq 0,
\]
whereas any gradient field integrates to zero around every closed loop. Different paths between the same endpoints give different answers, so no consistent potential can be defined.

\section{The Topological Twist}
\label{sec:topology}

Section~\ref{sec:gradient} established that $P_y = Q_x$ is \emph{necessary} for a potential to exist. Is it \emph{sufficient}? Here the subject takes its most surprising turn: the answer depends not on the field but on the \emph{shape of its domain}.

\begin{example}[The punctured plane]
Consider, on $\R^2 \setminus \{(0,0)\}$,
\[
(P,Q) = \left( \frac{-y}{x^2+y^2},\; \frac{x}{x^2+y^2} \right).
\]
A direct computation shows $P_y = Q_x$ at every point of the domain: both equal $(y^2 - x^2)/(x^2+y^2)^2$. Yet integrating around the unit circle again gives $2\pi \neq 0$, so no potential exists on the full punctured plane.
\end{example}

What went wrong? \emph{Locally}, nothing: near any point away from the origin, this field does have a potential, namely the polar angle $\theta = \arctan(y/x)$ (suitably interpreted in each half-plane). One can verify that $\nabla \theta$ is exactly the field above. But the polar angle cannot be defined as a continuous single-valued function on the whole punctured plane: traverse a loop around the origin and $\theta$ returns increased by $2\pi$. The candidate potential exists everywhere locally and fails to exist globally, and the value $2\pi$ of the loop integral is precisely the size of the mismatch. The obstruction lives in the hole.

The positive result delimiting this phenomenon is:

\begin{theorem}[Poincar\'e lemma, plane case]
Let $D \subseteq \R^2$ be open and \emph{simply connected} (connected, with no holes: every closed loop in $D$ can be continuously shrunk to a point within $D$). If $P, Q$ are continuously differentiable on $D$ and $P_y = Q_x$ throughout $D$, then $(P,Q) = \nabla f$ for some $f$ on $D$, unique up to a constant.
\end{theorem}

On a disk, a rectangle, or the whole plane, the compatibility condition is the whole story. On an annulus or a punctured plane, it is not: a field can pass the local test at every single point and still fail globally, and whether this can happen is determined by the topology of the domain alone.

This is a genuinely new phenomenon with no one-variable analogue---or rather, its one-variable shadow is the trivial remark from Section 2 that a disconnected domain permits one constant per piece. Connectivity governs uniqueness of the potential; \emph{simple} connectivity governs its existence. Local invertibility of the derivative is a matter of analysis; global invertibility is a matter of topology. The systematic measurement of the gap between ``locally invertible'' and ``globally invertible''---between \emph{closed} and \emph{exact}, in the language of the next section---is the starting point of de Rham cohomology, which converts questions about holes in spaces into linear algebra.

\section{Reunification: The Exterior Derivative and Stokes}

The preceding sections dismantled the FTC into fragments: one inverse theory for $\partial/\partial x$, another for the mixed partial, another for the gradient, each with its own antiderivative, kernel, and fundamental theorem. The final movement reassembles them.

The unifying device is the language of \emph{differential forms}. In the plane there are three kinds:
\begin{itemize}[itemsep=2pt]
\item \textbf{0-forms}: ordinary functions $f(x,y)$;
\item \textbf{1-forms}: expressions $\omega = P\dd x + Q\dd y$, the natural integrands of line integrals;
\item \textbf{2-forms}: expressions $h\dd x \wedge \dd y$, the natural integrands of double integrals.
\end{itemize}
A single operator, the \emph{exterior derivative} $d$, promotes each kind to the next:
\[
\{\text{0-forms}\} \xrightarrow{\;d\;} \{\text{1-forms}\} \xrightarrow{\;d\;} \{\text{2-forms}\},
\]
defined by
\[
df = \frac{\partial f}{\partial x}\dd x + \frac{\partial f}{\partial y}\dd y,
\qquad
d(P\dd x + Q\dd y) = \left( \frac{\partial Q}{\partial x} - \frac{\partial P}{\partial y} \right) \dd x \wedge \dd y.
\]
The first arrow is the gradient in thin disguise. The second packages exactly the combination $Q_x - P_y$ that governed everything in Sections~\ref{sec:gradient} and~\ref{sec:topology}. Two pieces of standard terminology: a form $\omega$ is \emph{closed} if $d\omega = 0$, and \emph{exact} if $\omega = d\alpha$ for some form $\alpha$ of one degree lower.

In this language the scattered facts snap together:
\begin{itemize}[itemsep=2pt]
\item Clairaut's theorem becomes the single identity $d(df) = 0$---``the derivative of a derivative is zero.'' Every exact form is closed.
\item The compatibility condition $P_y = Q_x$ says precisely that $\omega = P\dd x + Q\dd y$ is closed---the necessary condition for it to be exact, i.e., to be $df$ for a potential $f$.
\item The Poincar\'e lemma says: on a simply connected domain, closed implies exact. The punctured-plane example exhibits a closed form that is not exact, and de Rham cohomology is the systematic study of the quotient (closed forms modulo exact forms), which vanishes exactly when the domain has no holes of the relevant dimension.
\end{itemize}

And all the fundamental theorems collapse into one.

\begin{theorem}[Generalized Stokes theorem]
For a suitable oriented region $\Omega$ with boundary $\partial\Omega$, and a differential form $\omega$ of the appropriate degree,
\[
\int_{\partial\Omega} \omega = \int_{\Omega} d\omega.
\]
\end{theorem}

In words: \emph{integrating a derivative over a region equals evaluating the original object on the boundary}. Every fundamental theorem met above is a special case:
\begin{itemize}[itemsep=2pt]
\item Take $\Omega = [a,x]$, an interval, and $\omega = F$, a 0-form. The boundary of the interval is the two endpoints, the right one counted positively and the left negatively, and Stokes reads
\[
F(x) - F(a) = \int_a^x F'(t)\dd t,
\]
the one-variable FTC.
\item Take $\Omega = \gamma$, a curve from $A$ to $B$, and $\omega = f$. Stokes reads $f(B) - f(A) = \int_\gamma df$, the gradient theorem.
\item Take $\Omega$ a plane region and $\omega = P\dd x + Q\dd y$. Stokes reads
\[
\oint_{\partial\Omega} P\dd x + Q\dd y = \iint_\Omega (Q_x - P_y)\dd x \dd y,
\]
which is Green's theorem---and which explains, in one line, why closed fields have vanishing loop integrals around contractible loops.
\item In three dimensions the same statement specializes to the classical Stokes theorem and the divergence theorem.
\end{itemize}

The rhetorical center of the whole subject is this: the duality between differentiation and integration does not disappear in higher dimensions---it becomes a duality between \textbf{differentiation and taking boundaries}. The operator $d$ and the boundary operator $\partial$ are adjoint to one another under the integration pairing, and the identity $d \circ d = 0$ mirrors the geometric fact that a boundary has no boundary ($\partial \circ \partial = 0$: the boundary of a disk is a circle, and a circle has no endpoints). Seen from this height, the one-variable FTC was secretly a statement about boundaries all along; the boundary of $[a,x]$ just happened to be simple enough---two signed points---that no one needed to say so.

\section{Conclusion: The Fate of ``$+C$''}

The essay's protagonist has been the lost information, and its biography can now be written in full. In one variable it is a constant: the kernel of $d/dx$. Under a single partial derivative it inflates to an arbitrary function $g(y)$: the kernel of $\partial/\partial x$. Under the mixed partial it becomes a sum of two arbitrary one-variable functions, $g_1(x) + g_2(y)$---the kernel of $\partial^2/\partial x\,\partial y$, and the secret content of d'Alembert's traveling waves. Under the gradient it shrinks back to a single constant, but the inverse problem becomes conditionally solvable, with the compatibility condition $P_y = Q_x$ standing guard. And globally, on domains with holes, a new kind of lost information appears that is not the kernel of anything local: the cohomology class, measured by loop integrals such as the $2\pi$ of the punctured plane.

\begin{center}
\begin{tabular}{@{}llll@{}}
\toprule
Operator & Antiderivative & Lost information & Solvable? \\
\midrule
$d/dx$ & $\int_a^x \cdot \dd t$ & constant $C$ & always \\
$\partial/\partial x$ & $\int \cdot \dd x$, $y$ fixed & function $g(y)$ & always \\
$\partial^2/\partial x\,\partial y$ & $\iint \cdot \dd y\dd x$ & $g_1(x) + g_2(y)$ & always \\
$\nabla$ & line integral from base point & constant $C$ & iff $P_y = Q_x$ (+ topology) \\
$d$ (unified) & varies by degree & closed/exact gap & measured by cohomology \\
\bottomrule
\end{tabular}
\end{center}

Two morals. First, ``indefinite multiple integrals'' exist---the iterated partial antiderivative genuinely inverts the mixed partial---but textbooks pass over them because multivariable integration is taught in its measuring role rather than its inverting role, and because the inverting role is distributed among partial antiderivatives, potential functions, and, eventually, differential forms. Second, and more strikingly: the innocent question ``can I undo differentiation?''---pursued honestly in two variables---forces the invention of topology. The failure of integration to invert differentiation is not a defect of calculus; it is a detector of holes.

\appendix

\section{Exercises}

\begin{enumerate}[itemsep=4pt]
\item Solve $f_{xy} = 6xy$ by double indefinite integration, exhibit the full kernel, and verify by differentiation that the kernel terms vanish under $\partial^2/\partial x\,\partial y$.
\item Attempt to find $f$ with $\nabla f = (y, -x)$ by the recovery algorithm of Section~\ref{sec:gradient}, and identify the exact step at which the procedure fails. Confirm the failure by checking $P_y$ against $Q_x$.
\item Run the recovery algorithm on $\nabla f = (y e^{xy},\; x e^{xy} + 2y)$ and confirm $f = e^{xy} + y^2 + C$.
\item Verify by direct computation that the punctured-plane field $\left( \frac{-y}{x^2+y^2}, \frac{x}{x^2+y^2} \right)$ satisfies $P_y = Q_x$, then compute its integral around the unit circle (parametrize by $x = \cos\theta$, $y = \sin\theta$) and obtain $2\pi$.
\item Show that on the half-plane $x > 0$ the same field has the potential $f = \arctan(y/x)$, by differentiating. Why does this not contradict the nonexistence of a global potential?
\end{enumerate}

\section{Self-Check Questions}

Cover the answers; attempt each question before revealing it.

\begin{enumerate}[itemsep=8pt]

\item \textbf{The FTC consists of two statements. What are they, and which direction loses information?}\\
\emph{Answer:} (1) $\frac{d}{dx}\int_a^x f = f$: differentiating an integral recovers the integrand exactly. (2) $\int_a^x F' = F(x) - F(a)$: integrating a derivative recovers $F$ up to the constant $F(a)$. Only the second direction loses information.

\item \textbf{In the phrase ``lost information,'' what algebraic object does the lost information correspond to?}\\
\emph{Answer:} The kernel of the differential operator---the set of objects the operator sends to zero. For $d/dx$ on an interval, this is the constants.

\item \textbf{What is the kernel of $\partial/\partial x$ acting on functions of $(x,y)$, and what does this imply about ``constants of integration'' for partial antiderivatives?}\\
\emph{Answer:} All functions of $y$ alone. The constant of integration inflates into an arbitrary function $g(y)$.

\item \textbf{Do indefinite multiple integrals exist? What operator does the iterated indefinite integral $\iint \cdot \,dy\,dx$ invert?}\\
\emph{Answer:} Yes. It inverts the mixed partial $\partial^2/\partial x\,\partial y$, with lost information $g_1(x) + g_2(y)$.

\item \textbf{What is the kernel of $\partial^2/\partial x\,\partial y$, and how do you verify it in both directions?}\\
\emph{Answer:} $\{g_1(x) + g_2(y)\}$. Forward: each term is killed by one of the two partials. Converse: $f_{xy}=0$ means $f_x$ depends only on $x$; integrate to get $g_1(x) + g_2(y)$.

\item \textbf{How does d'Alembert's solution of the wave equation relate to this kernel?}\\
\emph{Answer:} In characteristic coordinates $\xi = x - ct$, $\eta = x + ct$, the wave equation becomes $u_{\xi\eta} = 0$, whose general solution is the kernel: $u = F(\xi) + G(\eta)$, two traveling waves.

\item \textbf{Why do textbooks rarely present double integrals as inverses of differentiation?}\\
\emph{Answer:} Multivariable integration is taught in its measuring role (volume, mass over regions); the inverting role is distributed among partial antiderivatives, potentials, and differential forms. But a double integral with variable upper limits does satisfy $\partial^2 F/\partial x\,\partial y = h$.

\item \textbf{Why is $\nabla f = (P,Q)$ a fundamentally different inverse problem from $f_{xy} = h$?}\\
\emph{Answer:} It is overdetermined: two equations, one unknown. Solvability requires the compatibility condition $P_y = Q_x$, whereas every continuous $h$ is somebody's mixed partial.

\item \textbf{Where does the condition $P_y = Q_x$ come from?}\\
\emph{Answer:} From Clairaut's theorem (equality of mixed partials): if $P = f_x$ and $Q = f_y$, then $P_y = f_{xy} = f_{yx} = Q_x$.

\item \textbf{What is the kernel of $\nabla$ on a connected domain? On a disconnected one?}\\
\emph{Answer:} Constants; one constant per connected component.

\item \textbf{In the recovery algorithm for a potential, what is the ``matching step,'' and what signals failure?}\\
\emph{Answer:} After integrating $P$ in $x$ to get $f = \int P\,dx + g(y)$, differentiate in $y$ and equate to $Q$; this must yield an equation for $g'(y)$ free of $x$. A surviving $x$ signals that no potential exists.

\item \textbf{Give a field that is not a gradient, and the one-line reason.}\\
\emph{Answer:} $(-y, x)$: $P_y = -1 \neq 1 = Q_x$. Gradients cannot swirl.

\item \textbf{State the gradient theorem. What is the multivariable analogue of ``antiderivative as definite integral with variable endpoint''?}\\
\emph{Answer:} $\int_\gamma \nabla f \cdot d\mathbf{r} = f(B) - f(A)$ for any curve from $A$ to $B$. The analogue: $f(x,y) = \int_{(x_0,y_0)}^{(x,y)} P\,dx + Q\,dy$ along any path, well-defined when integrals are path independent.

\item \textbf{Is $P_y = Q_x$ sufficient for a potential to exist? State the precise result.}\\
\emph{Answer:} Only on simply connected domains (Poincar\'e lemma). On domains with holes it can fail.

\item \textbf{Describe the standard counterexample and what its loop integral measures.}\\
\emph{Answer:} $\left(\frac{-y}{x^2+y^2}, \frac{x}{x^2+y^2}\right)$ on the punctured plane: closed everywhere, yet its integral around the origin is $2\pi$---the jump of the would-be potential $\theta$ (the polar angle) after one loop.

\item \textbf{Why does the local potential $\theta = \arctan(y/x)$ fail to be a global potential?}\\
\emph{Answer:} The polar angle cannot be defined continuously and single-valuedly on the whole punctured plane; it increases by $2\pi$ around any loop enclosing the origin.

\item \textbf{Define \emph{closed} and \emph{exact} for a 1-form. Which implies which, and what theorem is the implication?}\\
\emph{Answer:} Closed: $d\omega = 0$ (i.e., $P_y = Q_x$). Exact: $\omega = df$. Exact implies closed, by $d(df) = 0$, which is Clairaut's theorem in disguise. The converse holds only on simply connected domains.

\item \textbf{What single identity of the exterior derivative encodes Clairaut's theorem, and what geometric fact mirrors it?}\\
\emph{Answer:} $d \circ d = 0$. It mirrors $\partial \circ \partial = 0$: a boundary has no boundary.

\item \textbf{State the generalized Stokes theorem and explain how the 1D FTC is a special case.}\\
\emph{Answer:} $\int_{\partial\Omega} \omega = \int_\Omega d\omega$. For $\Omega = [a,x]$ and $\omega = F$, the boundary is the signed pair of endpoints, giving $F(x) - F(a) = \int_a^x F'$.

\item \textbf{Which classical theorems are special cases of generalized Stokes?}\\
\emph{Answer:} The 1D FTC, the gradient theorem, Green's theorem, the classical Stokes theorem, and the divergence theorem.

\item \textbf{Into what does the derivative--integral duality transform in higher dimensions?}\\
\emph{Answer:} A duality between differentiation and taking boundaries: $d$ on forms is paired with $\partial$ on regions.

\item \textbf{Trace the fate of ``$+C$'' across all the operators discussed.}\\
\emph{Answer:} Constant ($d/dx$) $\to$ arbitrary function $g(y)$ ($\partial/\partial x$) $\to$ $g_1(x)+g_2(y)$ (mixed partial) $\to$ constant again ($\nabla$, connected domain) $\to$ globally, a cohomology class measuring the closed/exact gap on domains with holes.

\item \textbf{Why does asking ``can I undo differentiation?'' lead to topology?}\\
\emph{Answer:} Local solvability is governed by analysis (the compatibility condition), but global solvability depends on the shape of the domain: closed-but-not-exact forms detect holes. Measuring this gap is de Rham cohomology.

\end{enumerate}

\end{document}
